%%%%%%%%%%%%%%%%%%%%%%%%%%%%%%%%%%%%%%%%%%%%%%%%%%%%%%%%%%%%%%%%%%
%%         mmm  mmmmmm mm   m        m    m mmmmmm  mmmm        %% 
%%       m"   " #      #"m  #        #    # #      "   "#       %%
%%       #   mm #mmmmm # #m #        #    # #mmmmm     m"       %%
%%       #    # #      #  # #  """   #    # #        m"         %%
%%        "mmm" #mmmmm #   ##        "mmmm" #mmmmm m#mmmm       %%
%%                                                              %%
%%                                                              %%
%%      Grundlagen der Elektrischen Netzwerke, UE               %%
%%      Gruppe 5, Team F                                        %%
%%      Authors: Severin Wolf, Maximilian Seidler.              %%
%%%%%%%%%%%%%%%%%%%%%%%%%%%%%%%%%%%%%%%%%%%%%%%%%%%%%%%%%%%%%%%%%%
\documentclass[a4paper]{article}
\usepackage{amsmath}
\usepackage[utf8]{inputenc}
\usepackage[T1]{fontenc}
\usepackage[english]{babel}
\usepackage{geometry}
\usepackage{graphicx}
\usepackage{tikz}
\usepackage{listings}
\geometry{a4paper,left=3cm,right=2cm, top=2cm, bottom=2cm} 
\usepackage[EFvoltages, european, straightvoltages]{circuitikz}

%tikz
\ctikzset{resistor = european}
\usetikzlibrary{decorations.pathreplacing}

%no paragraph indent
\setlength{\parindent}{0pt}

%for math, that does not fit
\renewcommand*{\arraystretch}{1.3}
\newcommand\scalemath[2]{\scalebox{#1}{\mbox{\ensuremath{\displaystyle #2}}}}

\newcommand\blfootnote[1]{%
	\begingroup
	\renewcommand\thefootnote{}\footnote{#1}%
	\addtocounter{footnote}{-1}%
	\endgroup
}

\begin{document}
\pagestyle{empty} \enlargethispage*{25cm}\samepage{
\vspace*{-3cm}
\begin{center}
\begin{minipage}[!h]{16cm}
\hspace*{0.2cm}
\includegraphics[width=3.3cm]{./Figures/igte_logo}
\begin{tabular}{p{8cm}}
\vspace{0.2cm}
\centering{
\Large Institute of Fundamentals and Theory in
 Electrical Engineering\\
Graz University of Technology\\
~\\}
\end{tabular}
\includegraphics[width=3.3cm]{./Figures/TUG_logo}
\end{minipage}
		\Large
		\textbf{Fundamentals of electrical circuits} \vspace*{0.5cm}\\
		\textbf{2. Homework}\\Th$\acute{\text{e}}$venin Equivalent Circuit and Linearity
		\vspace*{0.5cm}
		
		\large
		18 March 2021
	\end{center}}
	
	\vspace*{1cm}
	
	%%%%%%%%%%%%%%%%%%%%%%%%%%%%%%%%%%%%%%%%%%%%%%%%%%%%%%%%%%%%%%%%%%%%
	%%%%%%%%%%%%%%%%%%%%%%%%%%%%%%%%%%%%%%%%%%%%%%%%%%%%%%%%%%%%%%%%%%%%
	\section*{Assignment 1}
	Consider the following circuit. The voltage source $U_{S1}$ is a \textbf{current-controlled voltage source}.\\
	The Th$\acute{\text{e}}$venin equivalent should be determined for the circuit on the left hand side with respect to the terminals \textbf{(a,b)} by using the \textbf{\textcolor[rgb]{1,0,0}{modified} node voltage method}.
	\begin{enumerate}
		\item  Mark the currents and voltages across every element, the missing nodes and all node-voltages in the part of the circuit to be replaced. (0.5P)
		\item Find the Th$\acute{\text{e}}$venin equivalent circuit with respect to the terminals $\mathbf{a}$ and $\mathbf{b}$ representing the left hand side of the circuit ($U_{Th}$,$R_{Th}$). (3P)
		\begin{itemize} 
			\item Therefor apply the modified node voltage method to the part of the circuit to be replaced. The matrix should be determined by inspection. Explain your thoughts briefly in the protocol (equations are not necessary). Write down the additional conditions needed to solve the system of equations. Solve the system of equations in matrix form $\mathbf{A}\cdot \mathbf{x} = \mathbf{b}$ (include the results!) in Matlab and determine the open-circuit voltage $U_{\text{oc}}$ at the terminals $\mathbf{a}$ and $\mathbf{b}$.
			\item Apply a test voltage $U_T$ in order to calculate the Th$\acute{\text{e}}$venin resistance $R_{Th}$. \\
			\textbf{Hint1:} Each independent source can be deactivated in LTspice by simply setting its value to 0. No additional changes to the circuit are needed.
		\end{itemize}
		\item Confirm the calculated results for $U_{Th}$ and $R_{Th}$ in LTspice. Simulate both circuits. In order to confirm that the the Th$\acute{\text{e}}$venin equivalent circuit and the original circuit are equivalent to each other, make a simulation of the original circuit and the one with the calculated Th$\acute{\text{e}}$venin equivalent. Compare the results. Overall four simulations are needed. (1.5P)		
	\end{enumerate} 
			\textbf{2P Bonus:}
			\begin{itemize} 
				\item Draw the given circuit with the previously calculated Th$\acute{\text{e}}$venin equivalent circuit on the left hand side.
				\item Adjust the value of the current source $I_{S3,\text{new}}$ so that the voltage across the resistor $R_6$ equals $U_{R_6,\text{new}}=20V$ and confirm your result in LTspice.
			\end{itemize}
\vspace*{3.5cm}
	\subsection*{Values:}
	$R_1 = 10\Omega$ \qquad $R_2 = 5\Omega$ \qquad $R_3 = 8\Omega$ \qquad $R_4 = 12\Omega$ \qquad $R_5 = 6\Omega$ \qquad $R_6 = 20\Omega$ \\
	$U_{S1} = \alpha \cdot I_{R_3}$  \qquad $U_{S2} = 5V$ \qquad $I_{S3} = 0.8A$ \qquad  $\alpha = 4 \Omega = 4 \frac{V}{A}$ \\
\blfootnote{Deadline: 25.03.2021 - 9:00; \qquad presentation: group no. 2}

\newpage	
\begin{figure}[h!] \centering    
\begin{circuitikz}
      %controlled Us_1
      \draw (0,0) to[cvsource, v_=$U_{S1}$, color=blue]  (0,3);
      %Us_2
      \draw (6,3) to[V, v=$U_{s2}$, color=blue]         (3,3);
      %Is_3
      \draw (11,3) to[I, i=$I_{s3}$, color =red]        (11,6);
      %R_1
      \draw (3,0) to[R, *-, l=$R_1$]                    (6,0);
      %R_2
      \draw (3,0) to[R, -*, l=$R_2$]                    (3,3);
      %R_3
      \draw (3,3) to[R, -*, l=$R_3$, i=$I_{R3}$]        (3,6);
      %R_4
      \draw (0,3) to[R, *-, l=$R_4$]                    (3,3);
      %R_5
      \draw (6,3) to[R, *-*, l=$R_5$]                    (6,6);
      %R_6
      \draw (9,3) to[R, *-*, l=$R_6$]                    (9,6);
      %connections
      \draw (0,3) to[short] (0,6) to[short] (7.5,6);
      \draw (6,0) to[short] (6,3) to[short] (7.5,3);
      %clamps + connections
      \draw (7.5,6) to[short, o-] (11,6);
      \draw (7.5,3) to[short, o-] (11,3);
      %ref node
      \draw (0,0) to[short]                             (3,0);
      %grnd
      \draw (0,0) to (0,0) node[ground]{};
      %nodes
      \node[above]              (n_1) at (3,6) {$n_1$};
      \node[above, xshift=3mm]  (n_3) at (3,3) {$n_2$};
      \node[above, xshift=3mm]  (n_4) at (6,3) {$n_3$};
      \node[below]              (ref) at (3,0) {ref};
      \node[above]              (a) at (7.5,6) {a};
      \node[above]              (b) at (7.5,3) {b};
\end{circuitikz}
\caption{Given circuit for assignment 2}
\label{fig:circuit_assignment2}
\end{figure}

\begin{figure}[h!] \centering    
\begin{circuitikz}[scale=0.75]
%%
%%      Wenn du  Stromrichtungen ändern willst, dann musst du
%%      bei den to[short] nach dem i ein < oder > machen ;)
%%
      %controlled Us_1
      \draw (0,0) to[cvsource, v_=$U_{S1}$, color=blue]  (0,5);
      %Us_2
      \draw (10,5) to[V, v=$U_{s2}$, color=blue]         (5,5);
      %Is_3
      \draw (18,5) to[I, i=$I_{s3}$, color =red]        (18,10);
      %R_1
      \draw                                     (5,0)
      to[short, *-]                             (6,0)
      to[R, l=$R_1$]                            (9,0)
      to[short, i<=$I_{R1}$, color=red]         (10,0);
      %R_2
      \draw                                     (5,0) 
      to[short]                                 (5,1)
      to[R, l=$R_2$]                            (5,4)
      to[short, i<=$I_{R2}$, color=red]         (5,5);
      %R_3
      \draw                                     (5,5)
      to[short, *-]                             (5,6)
      to[R, l=$R_3$]                            (5,9)
      to[short, i=$I_{R3}$, color=red]          (5,10)
      to[short, -*]                             (5,10);
      %R_4
      \draw                                     (0,5)
      to[short, *-]                             (1,5)
      to[R, l=$R_4$]                            (4,5)
      to[short, i=$I_{R4}$, color=red]          (5,5);
      %R_5
      \draw                                     (10,5)
      to[short, *-]                             (10,6)
      to[R, l=$R_5$]                            (10,9)
      to[short, i=$I_{R5}$, color=red]          (10,10)
      to[short, -*]                             (10,10);
      %R_6
      \draw                                     (15,5)
      to[short, *-]                             (15,6) 
      to[R, l=$R_6$]                            (15,9)
      to[short, i=$I_{R6}$, color=red]          (15,10)
      to[short, -*]                             (15,10);
      %connections
      \draw (0,5) to[short] (0,10) to[short]    (12.5,10);
      \draw (10,0) to[short] (10,5) to[short]   (12.5,5);
      %clamps + connections
      \draw (12.5,10) to[short, o-] (18,10);
      \draw (12.5,5) to[short, o-] (18,5);
      %ref node
      \draw (0,0) to[short]                             (5,0);
      %grnd
      \draw (0,0) to (0,0) node[ground]{};
      %nodes
      \node[above]              (n_1) at (5,10) {$n_1$};
      \node[above, xshift=3mm]  (n_3) at (5,5) {$n_2$};
      \node[above, xshift=3mm]  (n_4) at (10,5) {$n_3$};
      \node[below]              (ref) at (5,0) {ref};
      \node[above]              (a) at (12.5,10) {a};
      \node[above]              (b) at (12.5,5) {b};
\end{circuitikz}
\caption{Network with curents and (node-) voltages}
\label{fig:circuit_labeled}
\end{figure}

\end{document}
